\documentclass[../LuanVan.tex]{subfiles}
\begin{document}

\section{Tổng quan, cấu trúc chức năng và kiến trúc hệ thống}
\label{section:4.1}

Chương này trình bày cách hệ thống được xây dựng từ yêu cầu nghiệp vụ thành các nhóm chức năng, thành phần phần mềm, luồng xử lý và giao diện sử dụng. Các hình chụp giao diện trong chương có mục đích chỉ ra nơi người quản lý thực hiện từng chức năng và dữ liệu nào được hiển thị. Dữ liệu thử nghiệm, tiêu chí đạt và bằng chứng cho kết quả vận hành được tách sang Chương 5.

Hệ thống được thiết kế cho một cửa hàng, trong đó dữ liệu hình ảnh tại các khu vực được liên kết với hồ sơ khách hàng, lần mua sắm và dữ liệu giao dịch. Về mặt chức năng, hệ thống được chia thành ba nhóm chính: tiếp nhận và xử lý quan sát; quản lý CRM và lần mua sắm; phân tích và báo cáo. Về mặt phần mềm, các nhóm chức năng này được phân bổ cho giao diện CRM, máy chủ nghiệp vụ, dịch vụ xử lý ảnh và cơ sở dữ liệu.

\subsection{Cấu trúc chức năng của hệ thống}
\label{subsection:4.1.1}

Nhóm tiếp nhận và xử lý quan sát chịu trách nhiệm kiểm tra dữ liệu đầu vào, phát hiện khuôn mặt, phân loại biểu cảm, tạo đặc trưng khuôn mặt và lưu kết quả xử lý. Nhóm CRM và lần mua sắm quản lý khách hàng, sự đồng ý sử dụng dữ liệu khuôn mặt, khu vực, sản phẩm, đơn hàng và chuỗi quan sát của từng lần mua sắm. Nhóm phân tích và báo cáo tổng hợp phân bố biểu cảm, biến thiên theo thời gian, thay đổi giữa các khu vực và chất lượng dữ liệu. Cấu trúc này được thể hiện tại Hình~\ref{fig:function_structure_ch4}.

\begin{figure}[H]
    \centering
    \resizebox{0.98\textwidth}{!}{%
    \begin{tikzpicture}[
        node distance=7mm and 11mm,
        root/.style={draw, rounded corners=2pt, align=center, minimum height=13mm, text width=3.2cm, font=\small\bfseries},
        group/.style={draw, rounded corners=2pt, align=center, minimum height=12mm, text width=3.6cm, font=\small\bfseries},
        item/.style={draw, align=left, minimum height=17mm, text width=5.2cm, font=\scriptsize},
        flow/.style={-{Latex[length=2mm]}, thick}
    ]
        \node[root] (system) {Hệ thống CRM và phân tích biểu cảm};
        \node[group, right=of system, yshift=23mm] (observation) {Tiếp nhận và xử lý quan sát};
        \node[group, right=of system] (crm) {CRM và lần mua sắm};
        \node[group, right=of system, yshift=-23mm] (report) {Phân tích và báo cáo};
        \node[item, right=of observation] (observation_items) {Kiểm tra dữ liệu đầu vào; phát hiện khuôn mặt; phân loại biểu cảm; tạo đặc trưng khuôn mặt; lưu kết quả và lịch sử xử lý};
        \node[item, right=of crm] (crm_items) {Quản lý khách hàng, sự đồng ý, khu vực, sản phẩm, đơn hàng; liên kết quan sát theo lần mua sắm};
        \node[item, right=of report] (report_items) {Phân bố biểu cảm; biến thiên theo thời gian; thay đổi giữa các khu vực; báo cáo chất lượng dữ liệu};
        \draw[flow] (system) -- (observation);
        \draw[flow] (system) -- (crm);
        \draw[flow] (system) -- (report);
        \draw[flow] (observation) -- (observation_items);
        \draw[flow] (crm) -- (crm_items);
        \draw[flow] (report) -- (report_items);
    \end{tikzpicture}%
    }
    \caption{Cấu trúc các nhóm chức năng của hệ thống}
    \label{fig:function_structure_ch4}
\end{figure}

Ba nhóm chức năng không hoạt động tách rời. Kết quả của nhóm xử lý quan sát là đầu vào để hình thành lần mua sắm; dữ liệu CRM và lần mua sắm tiếp tục cung cấp ngữ cảnh cho nhóm phân tích. Cách phân nhóm này cũng xác định cấu trúc trình bày của các Mục~\ref{section:4.2}, \ref{section:4.3} và \ref{section:4.4}.

\subsection{Tác nhân và các trường hợp sử dụng chính}
\label{subsection:4.1.2}

Hệ thống có ba tác nhân theo nghĩa chức năng. Người quản lý thao tác trên CRM để quản lý dữ liệu và xem kết quả. Nguồn thu nhận gửi khung hình cùng ngữ cảnh khu vực, thời gian và mã sự kiện. Quản trị hệ thống cấu hình tài khoản, khu vực, phiên bản mô hình và theo dõi trạng thái dịch vụ. Dịch vụ xử lý ảnh và cơ sở dữ liệu là thành phần bên trong, không phải tác nhân sử dụng trực tiếp.

\begin{table}[H]
    \centering
    \caption{Tác nhân và các trường hợp sử dụng chính}
    \label{tab:actors_use_cases_ch4}
    \resizebox{\textwidth}{!}{%
    \begin{tabular}{|>{\raggedright\arraybackslash}p{3.2cm}|>{\raggedright\arraybackslash}p{5.8cm}|>{\raggedright\arraybackslash}p{7.0cm}|}
        \hline
        \textbf{Tác nhân} & \textbf{Trường hợp sử dụng} & \textbf{Kết quả cần tạo ra} \\
        \hline
        Người quản lý & Quản lý khách hàng, dữ liệu khuôn mặt, sản phẩm, đơn hàng và khu vực & Dữ liệu CRM được tạo, cập nhật và tra cứu theo quyền của người dùng \\
        \hline
        Người quản lý & Theo dõi lần mua sắm; kiểm tra, xác nhận hoặc xử lý lại quan sát & Quan sát được liên kết đúng ngữ cảnh và mọi thay đổi đều có lịch sử \\
        \hline
        Người quản lý & Xem phân bố biểu cảm, biến thiên theo thời gian, thay đổi giữa các khu vực và chất lượng dữ liệu & Kết quả tổng hợp được tính từ các bản ghi đủ điều kiện và có thể lọc \\
        \hline
        Nguồn thu nhận & Gửi ảnh kèm khu vực, thời điểm và mã sự kiện & Một sự kiện thu nhận và từ không đến nhiều quan sát khuôn mặt được hình thành \\
        \hline
        Quản trị hệ thống & Quản lý tài khoản, cấu hình và trạng thái vận hành & Các dịch vụ được kiểm soát và cấu hình được truy vết \\
        \hline
    \end{tabular}}
\end{table}

Việc tách trường hợp sử dụng theo tác nhân giúp xác định rõ đầu vào và đầu ra của từng chức năng. Đặc biệt, một khung hình không đồng nhất với một khách hàng: nguồn thu nhận tạo một sự kiện ở cấp khung hình, còn hệ thống tạo một quan sát riêng cho mỗi khuôn mặt được phát hiện. Vì vậy, cảnh có nhiều người vẫn được tiếp nhận và không bị loại chỉ vì có nhiều hơn một khuôn mặt.

\subsection{Kiến trúc các thành phần}
\label{subsection:4.1.3}

Hệ thống áp dụng cách tổ chức phân tách theo trách nhiệm. Giao diện CRM đảm nhiệm tầng trình bày; máy chủ nghiệp vụ là điểm điều phối các quy tắc và dữ liệu; dịch vụ xử lý ảnh cung cấp kết quả suy luận; PostgreSQL đảm nhiệm lưu trữ lâu dài. Máy chủ nghiệp vụ nằm ở trung tâm kiến trúc và là thành phần duy nhất được phép đồng thời gọi dịch vụ xử lý ảnh, truy cập dữ liệu CRM và quyết định việc liên kết một quan sát với khách hàng hoặc lần mua sắm. Cách tổ chức này ngăn giao diện tự thực hiện quy tắc nghiệp vụ, đồng thời ngăn dịch vụ xử lý ảnh tự thay đổi hồ sơ khách hàng.

Hình~\ref{fig:component_architecture_ch4} thể hiện hướng phụ thuộc giữa các thành phần. Trình duyệt nhận giao diện do máy chủ web cung cấp và gửi yêu cầu nghiệp vụ đến máy chủ nghiệp vụ thông qua giao diện trao đổi dữ liệu thống nhất. Nguồn thu nhận ảnh cũng gửi dữ liệu vào máy chủ nghiệp vụ thay vì gọi trực tiếp dịch vụ xử lý ảnh. Máy chủ nghiệp vụ chuyển tạm thời nội dung ảnh đến dịch vụ xử lý ảnh qua kênh giao tiếp nội bộ, nhận kết quả có cấu trúc rồi mới thực hiện đối sánh, liên kết và lưu trữ. Dịch vụ xử lý ảnh và cơ sở dữ liệu không phụ thuộc vào giao diện CRM; dịch vụ xử lý ảnh cũng không truy cập trực tiếp cơ sở dữ liệu.

\begin{figure}[H]
    \centering
    \resizebox{0.98\textwidth}{!}{%
    \begin{tikzpicture}[
        node distance=11mm and 15mm,
        block/.style={draw, rounded corners=2pt, align=center, minimum height=13mm, text width=3.1cm, font=\small},
        store/.style={draw, cylinder, shape border rotate=90, aspect=0.25, align=center, minimum height=15mm, text width=3.0cm, font=\small},
        flow/.style={-{Latex[length=2mm]}, thick},
        bidirectional/.style={{Latex[length=2mm]}-{Latex[length=2mm]}, thick},
        label/.style={font=\scriptsize, align=center}
    ]
        \node[block] (browser) {Trình duyệt\\người quản lý};
        \node[block, right=of browser] (web) {Giao diện CRM\\React và Nginx};
        \node[block, right=of web] (api) {Máy chủ nghiệp vụ\\FastAPI};
        \node[block, below=of api] (source) {Nguồn thu nhận\\ảnh và ngữ cảnh};
        \node[block, right=of api, yshift=11mm] (vision) {Dịch vụ xử lý ảnh};
        \node[store, right=of api, yshift=-13mm] (database) {PostgreSQL\\và pgvector};
        \draw[bidirectional] (browser) -- node[label, above] {Giao diện\\người dùng} (web);
        \draw[bidirectional] (web) -- node[label, above] {Yêu cầu và\\kết quả} (api);
        \draw[flow] (source) -- node[label, right] {ảnh, khu vực,\\thời gian} (api);
        \draw[bidirectional] (api) -- node[label, above] {Giao tiếp nội bộ} (vision);
        \draw[bidirectional] (api) -- node[label, below] {Truy vấn và\\giao dịch} (database);
    \end{tikzpicture}%
    }
    \caption{Kiến trúc và hướng phụ thuộc giữa các thành phần}
    \label{fig:component_architecture_ch4}
\end{figure}

Bảng~\ref{tab:system_components} tổng hợp công nghệ và trách nhiệm của từng thành phần trong kiến trúc. Công nghệ được lựa chọn theo vai trò của thành phần, trong khi quy tắc nghiệp vụ vẫn được tập trung tại máy chủ để tránh phân tán quyết định trên nhiều lớp.

\begin{table}[H]
    \centering
    \caption{Phân bổ trách nhiệm giữa các thành phần phần mềm}
    \label{tab:system_components}
    \resizebox{\textwidth}{!}{%
    \begin{tabular}{|>{\raggedright\arraybackslash}p{2.5cm}|>{\raggedright\arraybackslash}p{4.2cm}|>{\raggedright\arraybackslash}p{8.0cm}|}
        \hline
        \textbf{Thành phần} & \textbf{Công nghệ} & \textbf{Trách nhiệm trong thiết kế} \\
        \hline
        Giao diện CRM & React, TypeScript, Ant Design, ECharts, Nginx & Thu nhận thao tác của người dùng, gọi API, quản lý trạng thái hiển thị và trình bày hồ sơ, lần mua sắm, đơn hàng cùng các báo cáo. \\
        \hline
        Máy chủ nghiệp vụ & FastAPI, SQLAlchemy & Xác thực, kiểm tra yêu cầu, điều phối xử lý ảnh, đối sánh khách hàng, thực hiện quy tắc lần mua sắm, quản lý giao dịch dữ liệu và cung cấp API. \\
        \hline
        Dịch vụ xử lý ảnh & FastAPI, OpenCV; RetinaFace--MobileNet0.25; mô hình Emotion qua DeepFace; ArcFace & Kiểm tra và giải mã ảnh, phát hiện vùng khuôn mặt, phân loại biểu cảm, tạo véc-tơ đặc trưng và trả kết quả cùng thông tin mô hình. \\
        \hline
        Cơ sở dữ liệu & PostgreSQL 16, pgvector, Alembic & Lưu sự kiện thu nhận, các quan sát khuôn mặt, dữ liệu CRM, véc-tơ mẫu, lần mua sắm, đơn hàng, nhật ký và lịch sử xử lý; duy trì khóa, quan hệ và phiên bản cấu trúc dữ liệu. \\
        \hline
    \end{tabular}}
\end{table}

\textbf{Giao diện CRM.} Giao diện được xây dựng bằng React và TypeScript. Sau bước biên dịch, Nginx phục vụ các tệp tĩnh của giao diện cho trình duyệt và chuyển tiếp các yêu cầu nghiệp vụ đến máy chủ API dưới cùng một địa chỉ truy cập. Nginx không xử lý nhận dạng hay quy tắc CRM; vai trò của nó là phân phối giao diện, định tuyến yêu cầu và tách máy chủ ứng dụng khỏi kết nối trực tiếp của trình duyệt. Thành phần giao diện tổ chức các biểu mẫu, bảng dữ liệu, chức năng tải ảnh và biểu đồ; mọi dữ liệu nghiệp vụ đều được trao đổi với máy chủ thông qua API. Giao diện chỉ trình bày kết quả do máy chủ trả về, không trực tiếp tính báo cáo chính thức, không quyết định quyền truy cập và không kết nối đến PostgreSQL. Việc giữ ranh giới này giúp các quy tắc vẫn có hiệu lực khi yêu cầu đến từ một nguồn khác ngoài giao diện web.

\textbf{Máy chủ nghiệp vụ.} FastAPI tiếp nhận yêu cầu từ giao diện và nguồn thu nhận, kiểm tra dữ liệu vào, xác thực người dùng và quản lý phiên làm việc. Thành phần này sử dụng SQLAlchemy để thao tác dữ liệu theo giao dịch, gọi dịch vụ xử lý ảnh với thời gian chờ cấu hình được, sau đó áp dụng các quy tắc đối sánh khách hàng, tạo hoặc cập nhật lần mua sắm và liên kết đơn hàng. Các truy vấn tổng hợp báo cáo cũng được thực hiện tại máy chủ để mọi giao diện sử dụng cùng một cách tính. Vì vậy, đây là nơi tập trung các quyết định nghiệp vụ của hệ thống.

\textbf{Dịch vụ xử lý ảnh.} Dịch vụ được xây dựng như một thành phần độc lập, nhận ảnh từ máy chủ nghiệp vụ và trả về trạng thái của khung hình cùng danh sách kết quả theo từng khuôn mặt. Mỗi phần tử trong danh sách gồm chỉ số khuôn mặt, khung bao, mức tin cậy phát hiện, kết quả phân loại biểu cảm và véc-tơ đặc trưng. Dịch vụ không nhận mã khách hàng, không đọc dữ liệu CRM và không tự kết luận danh tính. Sau khi trả kết quả, ảnh truy vấn và véc-tơ tạm thời không được dịch vụ lưu thành hồ sơ nghiệp vụ. Thiết kế này cho phép thay cấu hình mô hình mà không làm thay đổi các chức năng CRM.

\textbf{Cơ sở dữ liệu.} PostgreSQL lưu đồng thời dữ liệu quan hệ và véc-tơ 512 chiều của mẫu khuôn mặt thông qua kiểu dữ liệu do pgvector cung cấp. Dữ liệu ảnh được tổ chức thành hai cấp: sự kiện thu nhận đại diện cho khung hình, còn quan sát đại diện cho từng khuôn mặt trong khung hình đó. Mã sự kiện là duy nhất ở cấp sự kiện thu nhận; trong bảng quan sát, cặp mã sự kiện và chỉ số khuôn mặt ngăn tạo trùng cùng một kết quả. Các khóa ngoại bảo toàn quan hệ giữa sự kiện, quan sát, khách hàng, lần mua sắm và đơn hàng. Alembic được sử dụng để quản lý tuần tự các thay đổi cấu trúc dữ liệu. Khi hệ thống khởi động bằng Docker Compose, các thay đổi này được áp dụng hoàn tất trước khi máy chủ nghiệp vụ bắt đầu phục vụ yêu cầu.

Trong cấu hình đóng gói, PostgreSQL phải đạt trạng thái sẵn sàng trước khi tiến trình cập nhật cấu trúc chạy; máy chủ nghiệp vụ chỉ khởi động sau khi cập nhật hoàn tất và dịch vụ xử lý ảnh đã sẵn sàng. Trọng số của mô hình được gắn vào dịch vụ xử lý ảnh ở chế độ chỉ đọc, còn dữ liệu PostgreSQL được đặt trong vùng lưu trữ riêng. Hướng phụ thuộc tổng quát là giao diện $\rightarrow$ máy chủ nghiệp vụ $\rightarrow$ dịch vụ xử lý ảnh hoặc cơ sở dữ liệu; không có phụ thuộc ngược từ tầng dưới lên giao diện.

\subsection{Giao tiếp giữa các thành phần}
\label{subsection:4.1.4}

Giao tiếp trong hệ thống được tổ chức xoay quanh máy chủ nghiệp vụ. Giao diện CRM và nguồn thu nhận không truy cập trực tiếp dịch vụ xử lý ảnh hoặc cơ sở dữ liệu; ngược lại, dịch vụ xử lý ảnh cũng không chủ động gọi giao diện hay đọc dữ liệu CRM. Máy chủ nghiệp vụ đóng vai trò trung gian tiếp nhận yêu cầu, kiểm tra điều kiện thực hiện, điều phối các thành phần phía sau và chuẩn hóa kết quả trả về. Ranh giới này bảo đảm mỗi thành phần chỉ trao đổi những dữ liệu cần thiết cho trách nhiệm của mình.

Các luồng giao tiếp chính được tổng hợp trong Bảng~\ref{tab:component_communication}. Dữ liệu trao đổi được mô tả theo ý nghĩa nghiệp vụ thay vì theo đường dẫn API cụ thể, qua đó làm rõ quan hệ giữa các thành phần mà không phụ thuộc vào cách tổ chức mã nguồn.

\begin{table}[H]
    \centering
    \caption{Các luồng giao tiếp chính giữa những thành phần của hệ thống}
    \label{tab:component_communication}
    \resizebox{\textwidth}{!}{%
    \begin{tabular}{|>{\raggedright\arraybackslash}p{3.0cm}|>{\raggedright\arraybackslash}p{3.0cm}|>{\raggedright\arraybackslash}p{5.2cm}|>{\raggedright\arraybackslash}p{5.0cm}|}
        \hline
        \textbf{Thành phần gửi} & \textbf{Thành phần nhận} & \textbf{Dữ liệu trao đổi} & \textbf{Mục đích} \\
        \hline
        Giao diện CRM & Máy chủ nghiệp vụ & Thông tin đăng nhập, điều kiện tra cứu và dữ liệu do người quản lý nhập & Thực hiện các chức năng quản lý, xác nhận dữ liệu và xem báo cáo \\
        \hline
        Nguồn thu nhận & Máy chủ nghiệp vụ & Ảnh, mã sự kiện, khu vực và thời điểm quan sát & Tạo một sự kiện quan sát có ngữ cảnh và có khả năng kiểm tra trùng lặp \\
        \hline
        Máy chủ nghiệp vụ & Dịch vụ xử lý ảnh & Ảnh cần phân tích và yêu cầu xử lý & Phát hiện khuôn mặt, phân loại biểu cảm và tạo véc-tơ đặc trưng \\
        \hline
        Dịch vụ xử lý ảnh & Máy chủ nghiệp vụ & Trạng thái khung hình và danh sách khung bao, biểu cảm, véc-tơ đặc trưng của từng khuôn mặt & Cung cấp các kết quả kỹ thuật độc lập để máy chủ tạo quan sát và áp dụng quy tắc nghiệp vụ \\
        \hline
        Máy chủ nghiệp vụ & Cơ sở dữ liệu & Truy vấn, dữ liệu tạo mới hoặc cập nhật và thông tin nhật ký & Bảo toàn hồ sơ CRM, quan sát, lần mua sắm và lịch sử thay đổi \\
        \hline
    \end{tabular}}
\end{table}

Đối với luồng tạo bản ghi quan sát, quá trình giao tiếp diễn ra theo năm bước. Thứ nhất, nguồn thu nhận gửi ảnh kèm mã sự kiện, khu vực và thời điểm ghi nhận. Thứ hai, máy chủ kiểm tra tính đầy đủ của yêu cầu, sự tồn tại của khu vực, định dạng thời gian, giới hạn của ảnh và khả năng trùng sự kiện. Chỉ dữ liệu vượt qua bước kiểm tra này mới được chuyển sang dịch vụ xử lý ảnh. Thứ ba, dịch vụ xử lý ảnh giải mã ảnh, phát hiện vùng khuôn mặt và thực hiện hai nhánh phân loại biểu cảm, tạo véc-tơ đặc trưng trên vùng cắt hợp lệ. Trạng thái xử lý và kết quả được trả về trong cùng một lần giao tiếp để tránh truyền lại ảnh nguồn.

Thứ tư, máy chủ nghiệp vụ tạo một sự kiện thu nhận cho hình ảnh và đối sánh véc-tơ khuôn mặt với các mẫu đã đăng ký theo ngưỡng nhận dạng đang được cấu hình. Khi xác định được khách hàng, máy chủ tìm lần mua sắm đang hoạt động hoặc tạo lần mua sắm phù hợp với thời điểm quan sát. Thứ năm, kết quả được lưu thành một bản ghi quan sát gắn với khách hàng, lần mua sắm và phiên bản mô hình. Vì vậy, dịch vụ xử lý ảnh chỉ cung cấp kết quả suy luận, còn quyết định liên kết dữ liệu vẫn thuộc về máy chủ nghiệp vụ.

Đối với các thao tác trên giao diện CRM, trình duyệt gửi yêu cầu đến máy chủ và nhận lại dữ liệu đã được kiểm tra theo quyền của người dùng. Các thao tác tạo, sửa, xác nhận khách hàng hoặc xử lý lại quan sát đều được máy chủ kiểm tra trước khi ghi vào cơ sở dữ liệu. Với các chức năng báo cáo, giao diện chỉ gửi điều kiện lọc và trình bày kết quả; việc lựa chọn bản ghi hợp lệ, tổng hợp số liệu và tính các chỉ số được thực hiện tập trung tại máy chủ. Cách tổ chức này giúp cùng một quy tắc được áp dụng cho mọi nguồn truy cập và tránh sai khác giữa số liệu hiển thị trên các màn hình.

Luồng giao tiếp cũng phân biệt lỗi tiếp nhận với lỗi xử lý. Yêu cầu thiếu thông tin hoặc có ngữ cảnh không hợp lệ bị dừng trước khi gọi dịch vụ xử lý ảnh và được ghi nhận để phục vụ báo cáo chất lượng dữ liệu. Ảnh không có khuôn mặt hoặc không thể giải mã vẫn tạo trạng thái ở cấp sự kiện thu nhận nhưng không tạo quan sát giả. Nếu dịch vụ xử lý ảnh không phản hồi trong thời gian cấu hình, máy chủ ghi nhận lỗi mô hình và không tiếp tục đối sánh khách hàng. Mã sự kiện duy nhất cho phép nguồn thu nhận gửi lại cùng hình ảnh mà không tạo thêm sự kiện hoặc quan sát trùng lặp.

Xét theo phạm vi, giao tiếp bên ngoài bao gồm trao đổi giữa giao diện CRM hoặc nguồn thu nhận với máy chủ nghiệp vụ; giao tiếp nội bộ bao gồm trao đổi giữa máy chủ với dịch vụ xử lý ảnh và cơ sở dữ liệu. Sự phân tách này vừa giới hạn quyền truy cập vào dữ liệu nhạy cảm, vừa giữ cho thành phần xử lý ảnh độc lập với cấu trúc của CRM. Đồng thời, máy chủ nghiệp vụ trở thành điểm kiểm soát chung đối với xác thực, kiểm tra dữ liệu, ghi nhật ký và tính nhất quán của giao dịch.

\section{Xây dựng chức năng tạo bản ghi quan sát}
\label{section:4.2}

Luồng này là cầu nối giữa dữ liệu hình ảnh và dữ liệu CRM. Một lượt xử lý bắt đầu từ hình ảnh có ngữ cảnh, đi qua kiểm tra đầu vào và xử lý ảnh, sau đó được diễn giải thành sự kiện thu nhận và bản ghi quan sát gắn với khách hàng. Trình tự xây dựng gồm: (1) tiếp nhận và kiểm tra hình ảnh; (2) phát hiện vùng khuôn mặt; (3) xử lý biểu cảm và đặc trưng khuôn mặt; (4) đối sánh với kho mẫu; (5) lưu quan sát và cập nhật lần mua sắm. Trường hợp ảnh hỏng, không có khuôn mặt hoặc không đối sánh được với khách hàng đã đăng ký dừng ở cấp sự kiện thu nhận.

\subsection{Tiếp nhận và kiểm tra dữ liệu}
\label{subsection:4.2.1}

Chức năng tiếp nhận được xây dựng để nhận đồng thời thông tin mô tả sự kiện và tệp ảnh. Bốn thông tin được kiểm tra gồm mã sự kiện, mã khu vực, thời gian ghi nhận có múi giờ và dữ liệu ảnh. Mã sự kiện có ràng buộc duy nhất tại bảng sự kiện thu nhận; vì vậy, gửi lại cùng một mã trả về tập kết quả đã có thay vì tạo thêm dữ liệu. Khu vực phải tồn tại và đang hoạt động. Ảnh rỗng hoặc vượt quá dung lượng cấu hình bị từ chối trước khi gọi dịch vụ xử lý ảnh.

Các yêu cầu thiếu mã sự kiện, thiếu khu vực, có khu vực không hợp lệ, thiếu thời gian hoặc thiếu ảnh được ghi vào kho theo dõi lỗi tiếp nhận. Cách lưu này giúp báo cáo chất lượng thể hiện được lỗi xảy ra trước khi một sự kiện thu nhận được tạo. Những trạng thái ảnh chỉ xác định được sau khi tiếp nhận, như không phát hiện được khuôn mặt hoặc không thể giải mã, được lưu ở cấp sự kiện. Nếu kết quả không đạt ngưỡng đối sánh khách hàng, sự kiện vẫn được giữ để chống xử lý trùng nhưng không có bản ghi quan sát đi kèm.

Màn hình tại Hình~\ref{fig:areas_ch4} trình bày các khu vực đã được khai báo để nguồn thu nhận gửi kèm dữ liệu. Thứ tự khu vực cũng được lưu để phát hiện vị trí bị thiếu trong một lần mua sắm.

\begin{figure}[H]
    \centering
    \includegraphics[width=0.98\textwidth]{Hinhve/Chuong4/4_06_khu_vuc.png}
    \caption{Danh sách các khu vực được xác định trước trong cửa hàng}
    \label{fig:areas_ch4}
\end{figure}

\subsection{Phát hiện khuôn mặt và hai nhánh nhận dạng}
\label{subsection:4.2.2}

Ảnh hợp lệ được giải mã và đưa vào mô hình phát hiện khuôn mặt. Khi không có khuôn mặt, hệ thống lưu trạng thái tại sự kiện thu nhận và không tạo quan sát. Vùng khuôn mặt hợp lệ được đưa vào nhánh phân loại bảy biểu cảm và nhánh tạo véc-tơ đặc trưng khuôn mặt. Nếu vùng khuôn mặt không thể cắt, sự kiện được ghi nhận với trạng thái lỗi và dừng xử lý.

Hai nhánh được xây dựng độc lập. Lỗi của mô hình FER không làm mất véc-tơ khuôn mặt đã tạo được; lỗi của ArcFace không làm mất nhãn biểu cảm đã phân loại được. Trạng thái phân loại biểu cảm và trạng thái nhận dạng được lưu riêng để phân biệt kết quả hợp lệ với trường hợp xử lý không thành công. Nhãn \textit{Neutral} chỉ được lưu khi là kết quả của mô hình, không được dùng thay cho trường hợp xử lý lỗi.

CRM cung cấp màn hình đăng ký mẫu khuôn mặt và tìm khách hàng bằng ảnh như Hình~\ref{fig:face_management_ch4}. Chỉ khách hàng đã đồng ý sử dụng dữ liệu khuôn mặt mới được tạo mẫu tham chiếu.

\begin{figure}[H]
    \centering
    \includegraphics[width=0.98\textwidth]{Hinhve/Chuong4/4_14_du_lieu_khuon_mat.png}
    \caption{Chức năng đăng ký và tìm khách hàng bằng ảnh}
    \label{fig:face_management_ch4}
\end{figure}

\subsection{Lưu kết quả và xử lý lại}
\label{subsection:4.2.3}

Sự kiện thu nhận lưu mã sự kiện, khu vực, thời gian ghi nhận, thời gian máy chủ nhận và trạng thái xử lý ảnh. Khi đối sánh thành công với khách hàng đã đăng ký, bản ghi quan sát lưu khung bao, mức tin cậy phát hiện, nhãn biểu cảm, điểm dự đoán của bảy lớp, khoảng cách nhận dạng, các trạng thái xử lý và thông tin mô hình. Nếu không đạt ngưỡng đối sánh, hệ thống không tạo bản ghi quan sát.

Màn hình chi tiết tại Hình~\ref{fig:observation_detail_ch4} hiển thị riêng trạng thái ảnh, trạng thái phân loại, trạng thái nhận dạng và các mô hình đã dùng.

Chức năng xử lý lại yêu cầu người quản lý cung cấp ảnh mới. Trước khi cập nhật, hệ thống lưu toàn bộ nhãn, mức tin cậy, điểm đối sánh, trạng thái và phiên bản mô hình cũ thành một phiên bản lịch sử. Vì vậy, kết quả mới không xóa dấu vết của lần xử lý trước.

\begin{figure}[H]
    \centering
    \includegraphics[width=0.98\textwidth]{Hinhve/Chuong4/4_16_chi_tiet_quan_sat.png}
    \caption{Chi tiết quan sát và lịch sử thay đổi kết quả}
    \label{fig:observation_detail_ch4}
\end{figure}

\section{Xây dựng chức năng CRM và liên kết lần mua sắm}
\label{section:4.3}

\subsection{Quản lý khách hàng và dữ liệu mua hàng}
\label{subsection:4.3.1}

Hồ sơ khách hàng gồm mã, họ tên, số điện thoại, thư điện tử, ảnh đại diện, trạng thái và thông tin đồng ý sử dụng dữ liệu khuôn mặt. Danh sách hỗ trợ tìm theo mã, tên hoặc số điện thoại. Khi mở hồ sơ, CRM hiển thị thông tin khách hàng cùng lịch sử đơn hàng như Hình~\ref{fig:customer_history_ch4}.

\begin{figure}[H]
    \centering
    \includegraphics[width=0.98\textwidth]{Hinhve/Chuong4/4_03_ho_so_lich_su_mua_hang.png}
    \caption{Hồ sơ khách hàng và lịch sử mua hàng}
    \label{fig:customer_history_ch4}
\end{figure}

\subsection{Quản lý sản phẩm và đơn hàng}
\label{subsection:4.3.2}

Sản phẩm được tổ chức theo nhóm và có mã, tên, giá hiện tại cùng trạng thái. Khi tạo đơn hàng, mã, tên và đơn giá được sao chép vào từng dòng đơn. Do đó, thay đổi giá sản phẩm sau này không làm thay đổi giá trị của đơn hàng đã phát sinh.

Nếu người dùng không chỉ định mã lần mua sắm, máy chủ tìm lần mua sắm của cùng khách hàng có khoảng thời gian bao phủ thời điểm tạo đơn. Nếu người dùng chỉ định mã, máy chủ kiểm tra cả khách hàng và thời gian trước khi tạo liên kết. Đơn hàng không thỏa mãn điều kiện vẫn được lưu mà không liên kết với lần mua sắm.

\begin{figure}[H]
    \centering
    \includegraphics[width=0.98\textwidth]{Hinhve/Chuong4/4_04_danh_muc_san_pham.png}
    \caption{Danh mục sản phẩm đã được tạo trong CRM}
    \label{fig:products_ch4}
\end{figure}

\begin{figure}[H]
    \centering
    \includegraphics[width=0.98\textwidth]{Hinhve/Chuong4/4_05_lich_su_mua_hang.png}
    \caption{Danh sách đơn hàng và dữ liệu lịch sử mua hàng}
    \label{fig:orders_ch4}
\end{figure}

\subsection{Khởi tạo, cập nhật và kết thúc lần mua sắm}
\label{subsection:4.3.3}

Khi một quan sát xác định được khách hàng, máy chủ tìm lần mua sắm đang hoạt động của người đó. Quan sát mới được thêm vào lần hiện tại nếu thời gian cách quan sát gần nhất không quá 30 phút. Nếu vượt quá khoảng này, lần cũ được đóng do quá thời gian và một lần mới được tạo. Nguyên nhân kết thúc \texttt{TIMEOUT} được lưu để phục vụ tra cứu.

Các quan sát được đọc theo thời gian ghi nhận, không theo thời gian máy chủ nhận. Khi hai khu vực khác nhau có cùng thời điểm trong một lần mua sắm, cả hai bản ghi được đánh dấu xung đột và không được dùng để tạo cặp thay đổi. Nếu chuỗi có khu vực đầu và khu vực sau nhưng thiếu một khu vực ở giữa theo thứ tự cấu hình, giao diện hiển thị rõ tên khu vực bị thiếu và không tự tạo nhãn thay thế.

Hình~\ref{fig:shopping_tracking_ch4} cho thấy danh sách lần mua sắm, chuỗi khu vực, cảnh báo khu vực thiếu, phân bố nhãn trong lần đang xem và đơn hàng liên quan. Người dùng có thể chọn từng quan sát để mở thông tin chi tiết.

\begin{figure}[H]
    \centering
    \includegraphics[width=0.98\textwidth]{Hinhve/Chuong4/4_07_theo_doi_mua_sam.png}
    \caption{Theo dõi các quan sát trong một lần mua sắm}
    \label{fig:shopping_tracking_ch4}
\end{figure}

\section{Xây dựng chức năng phân tích và báo cáo}
\label{section:4.4}

\subsection{Phân bố biểu cảm theo khu vực và thời gian}
\label{subsection:4.4.1}

Báo cáo phân bố chỉ sử dụng các bản ghi của khách hàng đã đăng ký có ảnh và kết quả phân loại biểu cảm hợp lệ. Máy chủ đếm riêng bảy nhãn trong từng khu vực và tính tỷ lệ trên tổng số quan sát hợp lệ của đúng khu vực đó. Người dùng có thể lọc theo khu vực và khoảng thời gian.

Hình~\ref{fig:distribution_screen_ch4} trình bày đồng thời biểu đồ và bảng số liệu. Bảng trình bày cả số lượng và tỷ lệ để tránh diễn giải một tỷ lệ cao được tạo từ quá ít quan sát.

\begin{figure}[H]
    \centering
    \includegraphics[width=0.98\textwidth]{Hinhve/Chuong4/4_08_phan_bo_bieu_cam.png}
    \caption{Phân bố bảy nhãn biểu cảm tại từng khu vực}
    \label{fig:distribution_screen_ch4}
\end{figure}

Trang tổng quan còn có mục theo dõi biến thiên trong cùng một khu vực. Máy chủ nhóm các quan sát theo khoảng 5, 15, 30 hoặc 60 phút rồi đếm từng nhãn trong mỗi khoảng. Hình~\ref{fig:timeline_area_ch4} là kết quả hiển thị trong cùng một ngày và cùng một khu vực.

\begin{figure}[H]
    \centering
    \includegraphics[width=0.98\textwidth]{Hinhve/Chuong4/4_13_bien_thien_theo_thoi_gian_khu_vuc.png}
    \caption{Biến thiên số quan sát theo thời gian trong cùng khu vực}
    \label{fig:timeline_area_ch4}
\end{figure}

\subsection{Thay đổi biểu cảm giữa các khu vực}
\label{subsection:4.4.2}

Máy chủ sắp xếp các bản ghi hợp lệ của từng lần mua sắm theo thời gian. Các bản ghi liên tiếp tại cùng một khu vực được xem là một nhóm. Người dùng chọn một trong ba cách lấy bản ghi đại diện: bản ghi đầu tiên, bản ghi cuối cùng hoặc bản ghi có mức tin cậy cao nhất. Sau đó, hệ thống tạo cặp từ hai nhóm khu vực liên tiếp và đếm tổ hợp nhãn trước--sau.

Kết quả gồm số lần và tỷ lệ trong cặp khu vực đang xét. Bản ghi có xung đột thời gian bị loại khỏi bước tạo cặp. Hình~\ref{fig:change_screen_ch4} thể hiện cùng số liệu dưới dạng sơ đồ luồng và bảng. Báo cáo chỉ cho biết nhãn nào được ghi nhận trước và sau khi khách hàng xuất hiện tại hai khu vực liên tiếp. Chẳng hạn, một cặp \textit{Neutral}--\textit{Happy} được tính là một lần thay đổi nhãn, nhưng không được diễn giải trực tiếp là trải nghiệm của khách hàng đã được cải thiện. Hệ thống chưa có dữ liệu khảo sát mức độ hài lòng để kiểm chứng cách diễn giải đó và cũng không áp dụng ma trận chuyển tiếp Markov.

\begin{figure}[H]
    \centering
    \includegraphics[width=0.98\textwidth]{Hinhve/Chuong4/4_09_thay_doi_bieu_cam.png}
    \caption{Số lượng và tỷ lệ thay đổi nhãn giữa các khu vực liên tiếp}
    \label{fig:change_screen_ch4}
\end{figure}

\subsection{Phân cụm chuỗi biểu cảm}
\label{subsection:4.4.3}

Chức năng phân cụm được hiện thực thành mô-đun thuần \texttt{sequence\_analysis.py} và một nhóm giao diện API riêng. API chỉ lấy các quan sát thuộc đúng một nguồn và một mã lần chạy. Với nguồn \texttt{SIMULATOR}, điều kiện lọc dùng \texttt{simulation\_run\_id}; với nguồn \texttt{CAMERA}, điều kiện lọc dùng \texttt{experiment\_run\_id}. Quy tắc này ngăn dữ liệu mô phỏng, dữ liệu ảnh KDEF và các lần chạy khác nhau bị trộn vào cùng một ma trận.

Mô-đun thực hiện bốn bước. Thứ nhất, các quan sát được nhóm theo \texttt{visit\_id}, sắp xếp theo \texttt{observed\_at} và gom các khung liên tiếp tại cùng điểm chạm thành một trạng thái đại diện. Thứ hai, \texttt{SequenceData} và \texttt{get\_distance\_matrix} của Sequenzo tạo ma trận Optimal Matching với \texttt{sm=CONSTANT}, \texttt{indel=1} và \texttt{norm=none}. Ma trận phải vuông, đối xứng, hữu hạn, không âm và có đường chéo bằng không trước khi được sử dụng. Thứ ba, \texttt{k\_medoids\_range} chạy PAM trên phạm vi \(K\) hợp lệ; \texttt{observation\_silhouette} tính silhouette của từng chuỗi. Thứ tư, mã ánh xạ chỉ số do thư viện trả về sang \texttt{visit\_id}, medoid và chuỗi gốc.

Kết quả được lưu trong ba bảng. \texttt{sequence\_analysis\_runs} lưu nguồn, tham số, phiên bản thuật toán, danh sách phương án \(K\), \(K\) được chọn, ASW, cảnh báo và SHA-256 của ma trận khoảng cách. \texttt{sequence\_cluster\_summaries} lưu medoid, kích thước, tỷ lệ, silhouette trung bình và khoảng cách trung vị của mỗi cụm. \texttt{sequence\_cluster\_assignments} lưu quan hệ giữa visit và cụm, chuỗi đầu vào, metadata điểm chạm, khoảng cách tới medoid và silhouette. Việc lưu theo lần chạy cho phép giữ nhiều kết quả với cấu hình khác nhau mà không ghi đè dữ liệu quan sát.

Các giao diện chính gồm:
\begin{itemize}
    \item \texttt{POST /api/v1/sequence-analyses}: tạo một lần phân tích từ mã lần phát dữ liệu;
    \item \texttt{GET /api/v1/sequence-analyses}: lấy lịch sử và trạng thái các lần chạy;
    \item \texttt{GET /api/v1/sequence-analyses/\{id\}/clusters}: lấy medoid và thống kê cụm;
    \item \texttt{GET /api/v1/sequence-analyses/\{id\}/assignments}: lấy danh sách visit đã gán cụm;
    \item \texttt{PATCH /api/v1/sequence-analyses/\{id\}/clusters/\{cluster\_id\}}: lưu tên diễn giải do người phân tích đặt.
\end{itemize}

Trên giao diện, tab \textit{Phân cụm chuỗi} cho phép tạo hoặc chọn một lần phân tích đã lưu. Phần tổng quan hiển thị số chuỗi được dùng và bị loại, \(K\), ASW, đồ thị ASW theo \(K\) và bảng điều kiện kích thước cụm. Hình~\ref{fig:sequence_overview_ch4} là ảnh chụp trực tiếp của lần chạy đầu-cuối bằng ảnh KDEF qua nguồn \texttt{CAMERA}; mã nguồn dữ liệu và trạng thái \texttt{COMPLETED} xuất hiện ngay trên màn hình để truy vết.

\begin{figure}[H]
    \centering
    \includegraphics[width=0.98\textwidth]{Hinhve/Chuong4/4_20_tong_quan_phan_cum_chuoi.png}
    \caption{Tổng quan lần phân tích chuỗi và quá trình chọn số cụm}
    \label{fig:sequence_overview_ch4}
\end{figure}

Hình~\ref{fig:sequence_medoids_ch4} hiển thị đầy đủ các medoid. Mỗi thẻ là một chuỗi có thật trong dữ liệu, kèm số thành viên, tỷ lệ, silhouette trung bình và khoảng cách trung vị. Ô tên cụm không được điền tự động; người phân tích chỉ đặt tên sau khi xem chuỗi đại diện.

\begin{figure}[H]
    \centering
    \includegraphics[width=0.98\textwidth]{Hinhve/Chuong4/4_21_cum_va_medoid.png}
    \caption{Các cụm và chuỗi medoid được trả từ PAM}
    \label{fig:sequence_medoids_ch4}
\end{figure}

Bảng assignment cho phép lọc theo cụm và mở ngược lại lần mua sắm. Liên kết này bảo đảm một kết quả phân cụm có thể truy vết về bốn quan sát gốc thay vì chỉ tồn tại như số liệu tổng hợp. Hình~\ref{fig:sequence_assignments_ch4} và Hình~\ref{fig:sequence_trace_ch4} minh họa hai đầu của liên kết này. Giao diện cũng hiển thị cảnh báo rằng cụm là kiểu diễn biến tương tự, không tự động biểu thị mức hài lòng.

\begin{figure}[H]
    \centering
    \includegraphics[width=0.98\textwidth]{Hinhve/Chuong4/4_22_danh_sach_gan_cum.png}
    \caption{Danh sách lần mua sắm cùng cụm, khoảng cách và silhouette}
    \label{fig:sequence_assignments_ch4}
\end{figure}

\begin{figure}[H]
    \centering
    \includegraphics[width=0.98\textwidth]{Hinhve/Chuong4/4_23_truy_vet_cum_ve_hanh_trinh.png}
    \caption{Truy vết một assignment về bốn quan sát của lần mua sắm}
    \label{fig:sequence_trace_ch4}
\end{figure}

Tài liệu OpenAPI của máy chủ đang chạy được dùng để đối chiếu các giao diện đã mô tả với đường dẫn thực tế. Hình~\ref{fig:sequence_api_docs_ch4} chỉ được chụp sau khi trang tài liệu tải được nhóm endpoint \path{/api/v1/sequence-analyses}.

\begin{figure}[H]
    \centering
    \includegraphics[width=0.98\textwidth]{Hinhve/Chuong4/4_24_tai_lieu_api_phan_cum_chuoi.png}
    \caption{Các endpoint phân cụm chuỗi trong tài liệu OpenAPI đang chạy}
    \label{fig:sequence_api_docs_ch4}
\end{figure}

Thành phần phân tích chuỗi được nối vào đúng dữ liệu đã có của hệ thống. Hồ sơ khách hàng vẫn được tạo ở chức năng CRM và liên kết với mẫu khuôn mặt; quan sát nguồn Camera vẫn đi qua dịch vụ thị giác và được ghép vào lần mua sắm trước khi có thể tham gia phân cụm. Hình~\ref{fig:kdef_customers_ch4} cho thấy năm hồ sơ thử nghiệm cùng ảnh đăng ký thực được máy chủ web phục vụ. Hình~\ref{fig:kdef_journey_ch4} cho thấy một hồ sơ trong số đó được tái hiện tại bốn điểm chạm; nguồn của từng quan sát được lưu là \textit{Camera}, không phải \textit{Simulator}.

\begin{figure}[H]
    \centering
    \includegraphics[width=0.98\textwidth]{Hinhve/Chuong4/4_25_danh_sach_khach_kdef.png}
    \caption{Hồ sơ và ảnh đăng ký KDEF được nạp vào chức năng khách hàng}
    \label{fig:kdef_customers_ch4}
\end{figure}

\begin{figure}[H]
    \centering
    \includegraphics[width=0.98\textwidth]{Hinhve/Chuong4/4_27_hanh_trinh_kdef_bon_diem_cham.png}
    \caption{Hành trình bốn điểm chạm được tạo từ các quan sát nguồn Camera}
    \label{fig:kdef_journey_ch4}
\end{figure}

Hệ thống tích hợp được đóng gói bằng Docker Compose. Tại thời điểm chụp bằng chứng, PostgreSQL, dịch vụ thị giác, API, dịch vụ mô phỏng và giao diện web đều ở trạng thái hoạt động; tiến trình migration đã kết thúc với mã thoát 0. Ngoài trạng thái container, chương trình kiểm chứng gọi trực tiếp các endpoint sẵn sàng của API, Vision, Simulator và trang web, tất cả đều trả HTTP 200. Hình~\ref{fig:docker_status_ch4} được sinh trực tiếp từ kết quả \texttt{docker compose ps --all} và các phản hồi HTTP; dữ liệu máy đọc tương ứng được lưu cùng artifact của lần chạy KDEF.

\begin{figure}[H]
    \centering
    \includegraphics[width=0.98\textwidth]{Hinhve/Chuong4/4_29_trang_thai_he_thong_docker.png}
    \caption{Trạng thái Docker Compose và các endpoint tại thời điểm kiểm chứng}
    \label{fig:docker_status_ch4}
\end{figure}

\subsection{Báo cáo chất lượng dữ liệu}
\label{subsection:4.4.4}

Báo cáo chất lượng tổng hợp ba nhóm thông tin. Nhóm thứ nhất là tổ hợp trạng thái ảnh, biểu cảm và danh tính. Nhóm thứ hai là các yêu cầu đầu vào không đủ điều kiện tạo bản ghi. Nhóm thứ ba gồm bản ghi đến chậm, xung đột thời gian, số khu vực bị thiếu và số lần mua sắm đủ điều kiện phân tích. Báo cáo dùng cùng bộ lọc khu vực và thời gian với báo cáo phân bố.

Hình~\ref{fig:data_quality_ch4} cho thấy các chỉ số được trình bày trước biểu đồ và bảng trạng thái. Các chỉ số này được dùng để xác định phần dữ liệu nào đã tham gia báo cáo biểu cảm và phần nào bị loại.

\begin{figure}[H]
    \centering
    \includegraphics[width=0.98\textwidth]{Hinhve/Chuong4/4_10_chat_luong_du_lieu.png}
    \caption{Báo cáo chất lượng của dữ liệu đã tiếp nhận và liên kết}
    \label{fig:data_quality_ch4}
\end{figure}

Các chức năng và quy tắc xử lý trên xác định cấu trúc phần mềm của hệ thống. Trạng thái vận hành của hệ thống tích hợp, dữ liệu chạy thử, kết quả kiểm thử và các giới hạn thực nghiệm được trình bày riêng tại Chương 5.

\end{document}
